% paper/sections/12_verification.tex
%
% §12  NS 方程式の物理的整合性検証 (Physical Consistency Verification)
%
% 検証の 3 層構造：
%   第10章 — コンポーネント単体（各演算子が設計精度を満たすか）
%   第11章 — 物理的整合性（組み上げた NS ソルバが物理法則を正しく記述するか）
%   第12章 — 二相流ベンチマーク（公表参照解との定量比較）
% ═══════════════════════════════════════════════════════════════════════════════
\section{NS 方程式の物理的整合性検証}
\label{sec:verification}

本稿の数値検証は以下の 3 層構造で設計されている：
\begin{enumerate}[noitemsep]
  \item[\textbf{第~\ref{sec:numerical_integrity}章}]
    \textbf{コンポーネント単体} — 各演算子・ソルバが設計精度を満たすか（済）
  \item[\textbf{本章}]
    \textbf{物理的整合性} — 組み上げた NS ソルバが物理法則を正しく記述するか
  \item[\textbf{第~\ref{sec:validation}章}]
    \textbf{二相流ベンチマーク} — 公表参照解との定量比較
\end{enumerate}
第~\ref{sec:numerical_integrity}章では CCD 微分演算子の $\Ord{h^6}$ 精度，
PPE 欠陥補正の格子収束，DCCD フィルタの散逸特性など，
個々のコンポーネントが設計精度を有することを確認した．
しかし，各部品が正しくても，それらの結合が物理法則と整合する保証はない．
本章では，Predictor--PPE--Corrector パイプライン全体を単一の NS ソルバとして動作させ，
解析解が既知の物理問題で整合性を検証する．
なお，本章のテストでは\textbf{成功と同時に限界も明らかになる}：
単相 NS ソルバは設計精度を完全に達成し（§~\ref{sec:ns_time_accuracy}），
低密度比（$\rho_l/\rho_g \leq 5$）の二相流は静止界面で安定動作する
（§~\ref{sec:advection_pressure_coupling}）一方，
§~\ref{sec:interface_crossing} の定量分析により
smoothed Heaviside 一括解法が $\rho_l/\rho_g \geq 10$ で破綻することを示す．
また §~\ref{sec:galilean_invariance} で移動界面＋表面張力（$\sigma > 0$）では
IPC が 2 ステップで発散することを実証し，
分相 PPE 解法による解決策を提示する．

\noindent\textbf{本章のスコープ：}
検証は物理的複雑さを段階的に引き上げ，以下の 6 段階で構成される：
\begin{enumerate}[noitemsep]
  \item \textbf{§~\ref{sec:force_balance}（力学的バランス）：}
        静止流体における重力・表面張力の釣合．
  \item \textbf{§~\ref{sec:ns_conservation}（保存則）：}
        運動エネルギー減衰と非圧縮条件の維持．
  \item \textbf{§~\ref{sec:ns_time_accuracy}（総合精度）：}
        NS 方程式の時間・空間収束次数と高 Re 安定性．
  \item \textbf{§~\ref{sec:advection_pressure_coupling}（二相流結合）：}
        CLS↔PPE 結合の物理的整合性（$\rho_l/\rho_g \leq 5$），
        HFE の適用範囲検証，Galilean 不変性．
  \item \textbf{§~\ref{sec:interface_crossing}（一括解法の適用限界）：}
        smoothed Heaviside の破綻閾値と分相 PPE 解法の必要性の数値的根拠．
  \item \textbf{§~\ref{sec:error_budget}（精度バジェット）：}
        全コンポーネントの実測精度を統合し，律速段階と今後の展望を示す．
\end{enumerate}

\noindent
各テストは理論解・解析解との定量比較により合否を判定する．
結果を踏まえた精度バジェットを §~\ref{sec:error_budget} に総括し，
第~\ref{sec:validation}章の二相流ベンチマークへと接続する．

\begin{table}[htbp]
\centering
\caption{物理的整合性の検証指標と許容基準}
\label{tab:verification}
\renewcommand{\arraystretch}{1.3}
\begin{tabular}{p{32mm}p{34mm}p{26mm}p{20mm}l}
\toprule
検証項目 & 理論値/解析解 & チェック方法 & 許容基準 & 節 \\
\midrule
静水圧バランス & $p = \rho g z$，$\bu = \bm{0}$ & $\|\bu\|_\infty$ & $< 10^{-5}$\textsuperscript{$\dagger$} & §\ref{sec:force_balance}\\
Laplace 圧均衡 & $\Delta p = \sigma/R$ & $\Delta p$ 相対誤差 & $< 1\%$ & §\ref{sec:force_balance}\\
運動エネルギー減衰 & $E_k(t) = E_k(0)\,e^{-4\nu t}$ & 相対誤差 & $< 10^{-4}$ & §\ref{sec:ns_conservation}\\
非圧縮条件 & $\bnabla\cdot\bu \equiv 0$ & $\|\bnabla\cdot\bu\|_\infty$ & $< 10^{-10}$ & §\ref{sec:ns_conservation}\\
NS 時間精度 & $\Ord{\Delta t^2}$（AB2+IPC） & $\Delta t$ 精緻化 & スロープ確認 & §\ref{sec:ns_time_accuracy}\\
NS 空間精度 & Kovasznay 解析解 & 格子精緻化 & スロープ確認 & §\ref{sec:ns_time_accuracy}\\
DCCD 非侵襲性 & $E_k^{\text{CCD}} \approx E_k^{\text{DCCD}}$ & 相対差 & $< 10^{-3}$ & §\ref{sec:ns_time_accuracy}\\
静止液滴安定性 & 寄生流れ有界 & 200 ステップ & 成長比 $< 5$ & §\ref{sec:advection_pressure_coupling}\\
HFE 適合性 & smoothed $p$: HFE 不要 & ON/OFF 比較 & 適用範囲の特定 & §\ref{sec:advection_pressure_coupling}\\
Galilean 不変性 & $\|\bu - \bu_0\|_\infty = 0$ & 1 周期後 & $< 10^{-10}$ & §\ref{sec:advection_pressure_coupling}\\
PPE 破綻限界 & DC 発散の閾値 & 残差履歴 & 閾値の特定 & §\ref{sec:interface_crossing}\\
分相 PPE 精度回復 & 密度比非依存 & 格子収束 & DC $k{=}3$: $\Ord{h^7}$ & §\ref{sec:interface_crossing}\\
\bottomrule
\multicolumn{5}{l}{\footnotesize \textsuperscript{$\dagger$}$N{=}32$ は CCD 壁面境界スキーム（$\Ord{h^5}$）の影響で $\Ord{10^{-6}}$ に留まる；$N{\geq}64$ では $\Ord{10^{-10}}$ 以下（§~\ref{sec:force_balance}）．} \\
\end{tabular}
\end{table}
